\documentclass{article}

\usepackage[letterpaper,margin=2.5cm]{geometry}
\usepackage{parskip}

\begin{document}

\section{Project Overview}

% (approximately 1 page)
% A paragraph or two describing what your application does, the category to which it belongs,
% and its key features. The project description should also describe the input and output to
% the program and how users interact with it (if it is interactive).

Our project belongs to ``Category 1 --- Multi-Process Application Using Pipes.''

[Project Name Here] is an advanced, cutting-edge, last-gen spinning cube and cuboid\footnote{also known as a rectangular prism} raytraced ASCII-art renderer.

The user can specify the 3-dimensional coordinates of each of the six vertices of zero or more cuboids as input to the program on startup.

While the program is running, the user will be find themselves in a three dimensional scene, where the cubes are rendered using ASCII characters to create an illusion of brightness differentials.
The user may walk around the scene to inspect the cube, provided they know C and can write that part for themself.

\section{Build Instructions}

% While we should be able to compile your project by running make with no arguments, there
% may be other requirements needed to build and run your program. List the exact commands
% needed to run your project, including any expected command-line arguments. If the project
% requires input files, include these files.

\section{Architecture Diagram}

% (no more than 1 page)
% A diagram (can be hand-drawn and photographed or produced with any drawing tool)
% showing:
% • All processes (or server/clients) and how they relate to each other.
% • The pipes or sockets that connect them, labeled with the type of data that flows
% through each connection.
% Depending on the project, this diagram may look quite different. The purpose of the diagram
% is to illustrate the overall structure of the program.

\section{Communication Protocol}


% For every message type, specify:
% • Sender and receiver
% • Encoding
% • Semantics and any invariants the receiver must uphold. In other words, how does
% the receiver know how many bytes to read, where to store the message, and what the
% message means.
% • Error handling
% See 2.1.4 for more details and an example.

\section{Concurrency Model}

% Be careful that your project avoids an overly complex concurrency model. You can rely on
% the synchronization provided by the message passing mechanisms (reading and writing on
% pipes or sockets).
% This section of the document describes how your application achieves concurrency:
% • For Category 1: explain when the parent forks, how it knows when workers are ready,
% and how it collects child processes.
% • For Category 2: explain your select loop, what is in the fd set, and how you avoid
% blocking on a single client.

\section{Error Handling and Robustness}

% Document at least three “bad behaviours” that could occur at runtime (e.g., a client discon-
% necting mid-message, a worker writing to a closed pipe, a malformed command) and explain
% how your code handles each.

\section{Team Contributions}

% One brief paragraph per team member describing their specific contributions.

\end{document}